%% kapa: exam=1 round=32 year=2021 period=2 form=A subject=appraisal_law subject_ko=감정평가관계법규 seq=19 number=19
%% source: Q-Net 감정평가사 기출문제 (cst00309) artlSeq=5208397 file=extracted/1008/5208397/A형/2021년도 제32회 감정평가사 1차 2교시 A형.pdf
%% conversion: kordoc 4.24.51 -> tools/gen1.py (auto); review=proofread-sonnet; flags=none
\begin{question}{19}
감정평가 및 감정평가사에 관한 법령상 감정평가법인에 관한 설명으로 옳지 않은 것은?
\begin{choices}
  \item 감정평가법인은 토지등의 이용 및 개발 등에 대한 조언이나 정보 등의 제공을 행한다.
  \item 감정평가법인은 토지등의 매매업을 직접 하여서는 아니 된다.
  \item 감정평가법인이 합병으로 해산한 때에는 이를 국토교통부장관에게 신고하여야 한다.
  \item 국토교통부장관은 감정평가법인이 업무정지처분 기간 중에 감정평가업무를 한 경우에는 그 설립인가를 취소할 수 있다.
  \item 감정평가법인의 자본금은 2억원 이상이어야 한다.
\end{choices}
\end{question}
\answer{4}
